\chapter{Reproducibility Artifacts}\label{appendix:artifacts}

\section{Purpose}

This appendix records the main components that support the AI-Scientist thesis. It is intended as a navigation aid
for examiners and future maintainers rather than as a replacement for the full project repository. The system is
implemented as a single Python package, \texttt{ai\_scientist}, with a Streamlit user interface and a FastAPI
service built on top of it.

\section{Core Modules}

\begin{table}[!htbp]
\centering
\caption{Core implementation modules of the AI-Scientist package.}
\label{tab:module-map}
\begin{tabular}{p{0.34\textwidth}p{0.55\textwidth}}
\toprule
Module & Responsibility \\
\midrule
\path{agents/orchestrator.py} & End-to-end workflow, early exits, and the re-verification loop. \\
\path{agents/retrieval.py} & Layered, source-prioritized retrieval and ranking. \\
\path{agents/live_retrieval.py} & Live retrieval from PubMed Central and arXiv. \\
\path{agents/claim_extraction.py} & Decomposition of abstracts into structured claims. \\
\path{agents/verification.py} & Evidence-bound verdicts and confidence computation. \\
\path{agents/critic.py} & Weakness critique and re-verification decisions. \\
\path{agents/synthesis.py} & Composition of the grounded executive summary. \\
\path{agents/report.py} & Structured and Markdown report assembly. \\
\path{models.py} & Typed dataclasses for papers, claims, evidence, and reports. \\
\path{config.py} & Settings: thresholds, source bonuses, and model configuration. \\
\path{utils.py} & Tokenization, Jaccard overlap, and query-coverage measures. \\
\path{scope.py} & Domain-general scope guard. \\
\path{llm.py} & Optional, provider-tolerant language-model refinement layer. \\
\path{corpus.py}, \path{ingestion.py} & Curated-corpus access and ingestion of user papers. \\
\path{baselines.py} & Single-agent and RAG comparison systems. \\
\bottomrule
\end{tabular}
\end{table}

\section{Interfaces}

\begin{table}[!htbp]
\centering
\caption{User-facing entry points.}
\label{tab:interface-map}
\begin{tabular}{p{0.30\textwidth}p{0.59\textwidth}}
\toprule
Entry point & Role \\
\midrule
\path{app.py} & Streamlit interface for uploading papers, selecting a model, asking questions, and inspecting per-claim verdicts, confidence, evidence, and the orchestration trace. \\
\path{ai_scientist/api.py} & FastAPI service exposing the same pipeline programmatically. \\
\bottomrule
\end{tabular}
\end{table}

\section{Verdict and Confidence Reference}

\begin{table}[!htbp]
\centering
\caption{Verdict space used by the Verification agent.}
\label{tab:verdict-map}
\begin{tabular}{p{0.26\textwidth}p{0.63\textwidth}}
\toprule
Verdict & Meaning \\
\midrule
\texttt{supported} & Retrieved evidence aligns with the claim; confidence scales with corroboration and source independence. \\
\texttt{contradicted} & Strong, multi-source counter-evidence dominates; the disagreement is reported explicitly. \\
\texttt{insufficient\_evidence} & No sufficiently aligned evidence, or the claim is overgeneralized; the system abstains. \\
\texttt{out\_of\_scope} & The query is not a research question and is declined before analysis. \\
\bottomrule
\end{tabular}
\end{table}

\section{Running the System}

The package is run with the source directory on the Python path. The deterministic pipeline operates fully offline
using the curated corpus and any uploaded papers; live retrieval from PubMed Central and arXiv is enabled through
the environment, optionally with an NCBI API key to raise the request rate; and the optional language-model
refinement is enabled when a provider API key is configured. Because the language model and live retrieval are
additive, the system produces the same structured assessment for a given question and evidence set whenever the
deterministic core alone is used.
